\scanpage{25}

$X$ espace topologique \deleted{compact} $G_1$ [triangulable compact de dim $\leq1$]

$S$ surface \uuline{orientée}

Nous voulons étudier les classes d'isotopie des plongements (modérés\ldots.) $X\xrightarrow{\varphi}S$.

Pour tout tel plongement\inserted{$\varphi$}, et pour tout sommet \deleted{de} \inserted{topologique} non isolé de $X$, il y a un \uline{ordre circulaire}\inserted{$\omega_x(\varphi)$} sur l'ens des arêtes (orientées) issues de $x$ induit par $\varphi$ \deleted{de} [donnée vide si le sommet $x$ est un bout, i.e. s'il n'y a qu'une seule arête issue de $x$ ; sinon, il y \deleted{a \unsure{un entier} \illegible{1}} a $\unsure{\nu|v}_x\geq3$ telles arêtes, et \deleted{le nombre de} l'ens des ordres circulaires possibles est un ens. $\Omega_x(X)$ à $(\nu_x-1)!$ éléments).

On trouve ainsi un premier invariant
\[
 \omega(\varphi)=(\omega_x(\varphi))_{x\in X'_0}\in\prod_{x\in X'_0}\Omega_x(X)
\]
\inserted{ens. des sommets de $X$ qui ne sont ni isolés ni des bouts}

On a une appl.
\[
 \mathrm{Pl}_{\mathrm{isot}.}(X,S)\xrightarrow{\omega}
 \prod_{x\in\unsure{x|X}'_0}\Omega_x(X)\overset{\text{déf}}{=}\Omega(X)
\]
et notre propos est d'étudier les fibres \inserted{$\mathrm{Pl}_{\mathrm{isot}}^\alpha(X,S)$} de $\omega$, plus précisément, pour $X$ et un $\alpha=(\alpha_x)_{x\in X'_0}\in\deleted{\unsure{\prod\Omega_x}}\Omega(X)$ fixés, d'étudier pour tt $S$ les plongements $\varphi:X\hookrightarrow S$ tels que $\omega(\varphi)=\alpha$.

Pour ceci, \deleted{\unsure{notons le}} \deleted{D}définissons pour tt \unsure{plongement} ``le'' voisinage tubulaire \inserted{fermé} relatif à celui-ci. On trouve

\uline{Th 1} Pour $\alpha$ fixé, ce voisinage tubulaire \deleted{est \unsure{défini} \illegible{1}}\inserted{ne dépend}, à homéomorphismes près induisant l'identité $X$ [de façon \unsure{lui-même|unique} à isotopie près] que de $\alpha$, soit $\unsure{U|\mathcal U}_\alpha$, et non de $S$. \deleted{\unsure{L'on|Donc}} On trouve alors une application
\[
 \mathrm{Pl}_{\mathrm{isot}}^\alpha(X,S)\longrightarrow\mathrm{Pl}_{\mathrm{isot}}(U_\alpha,S)
\]
Celle-ci est une bijection
% archive sheet number: 24
